\begin{pastexam}[label={px:nccu-114-dsa-09}]{政大資科所 114 學年度 第 9 題}
（單選；每題 4 分）

(4\%) Given an undirected connected graph $G$ of $n$ vertices and $m$ edges, which of the following
statement concerning spanning trees is \textbf{false}?
\begin{enumerate}[label=\Alph*.]
\item The minimum spanning tree should not include the heaviest edge in any cycle of $G$.
\item The maximum spanning tree can be found by modifying Kruskal's algorithm to sort edges in
descending order by weight.
\item Each spanning tree contains exactly $n$-1 edges.
\item Both Prim's and Kruskal's algorithms are greedy algorithms.
\item Heap is commonly used to implement Kruskal's Algorithm for finding minimum spanning tree
efficiently.
\end{enumerate}
\end{pastexam}

\begin{pastsoln}
答案是 E。 這題要注意 A 在權重有相同時說得不夠嚴謹，見下面 A 的說明；整體來看，題目要考的錯誤
敘述是 E。

逐項檢查：
\begin{itemize}
\item A.~這是環性質（資結書定理~24.3）：環上\textbf{唯一}最重的邊不在任何一棵最小生成樹裡。 如果
環上有好幾條邊一樣重，例如三角形三條邊都是 $1$，任何一棵最小生成樹都含其中兩條「最重的邊」，這時
A 字面上不成立；定理在這種情形只保證「有一棵」最小生成樹不含它。 在一般出題時「權重互不相同」的
默認下，A 是對的。
\item B.~對。 Kruskal 演算法的正確性只靠「依序考慮、不形成環就加入」這個貪心選擇；把邊由大到小
排序，同樣的證明（把權重全部取負號）給出最大生成樹。
\item C.~對。 $n$ 個頂點的樹恰好有 $n - 1$ 條邊，生成樹是包含全部頂點的樹。
\item D.~對。 兩者每一步都選目前最便宜、而且不破壞條件的邊，選了就不再反悔，正確性都由切割性質
（資結書定理~24.1）保證，是貪心法。
\item E.~錯，這是答案。 Kruskal 演算法（資結書虛擬碼~24.2）的兩個關鍵是\textbf{把邊排序}和
用\textbf{互斥集合}（union--find）判斷兩端是否已經連通；瓶頸是排序 $O(m \log m)$，判斷環用互斥集合
幾乎是常數時間（資結書第~24.2~節的分析）。 堆積才是 Prim 演算法的核心資料結構（資結書虛擬碼~24.3，
用最小堆積取出 $\mathit{key}$ 最小的頂點）。 用最小堆積代替排序、一條一條取出最小的邊也可以，
但那只是排序的一種做法，不是讓 Kruskal 有效率的關鍵，說它「常用堆積實作」是把 Kruskal 和 Prim
混在一起了。
\end{itemize}
\end{pastsoln}
